\section{Scope, limits, and out-of-scope claims}\label{sec:scope-limits}

Sections~\ref{sec:primitive-roles}--\ref{sec:conditional-lift} now contain the
definitions, theorem ladder, and conditional arithmetic lift that make up the
proof-carrying core of the manuscript. This final section records the paper's
perimeter: its continuity with the earlier Six-Birds line, the reason it is
theorem-led conditional rather than hybrid-led, and the stronger claims it
explicitly does not make.

\subsection*{Continuity with the earlier Six-Birds line}

The structural core of the paper is continuous with the earlier Six-Birds work
on fixed closure, packaging, and selective primitive activation. The
fixed-package saturation perspective comes from the structural closure line of
\emph{Six Birds: Foundations of Emergence Calculus}, where fixed closure does not
itself generate open-ended novelty~\cite{Tsiokos2026Foundations}. The treatment
of mathematics and formal systems as packaged, theory-relative objects is
continuous with \emph{To Count a Stone with Six Birds: A Mathematics is A
Theory}~\cite{Tsiokos2026ToCount}.

The manuscript also inherits a caution from the PICA line: richer primitive
interaction is real and should not be flattened into a single proxy merely
because one is pursuing a theorem route~\cite{Tsiokos2026ToCreate}. At the same
time, the current paper does not attempt to reproduce the full interactional
scope of PICA. It isolates one theorem-bearing slice from that broader setting.

Two earlier papers are especially important for the present role cleanup.
\emph{To Chart a Stone with Six Birds} shows that different projects may occupy
different regions of the primitive/control space rather than activating all six
primitives symmetrically at once~\cite{Tsiokos2026ToChart}. \emph{To Classify a
Stone with Six Birds} sharpens that point by separating class-defining
primitives from parameterizing or diagnostic ones~\cite{Tsiokos2026ToClassify}.
The present paper adopts that style of specialization for foundations: \(P5\) is
the fixed-package object, \(P1/P2\) are the package-change mechanisms, \(P6\) is
the frozen evaluation regime, \(P4\) is the support/index axis, and \(P3\) is
diagnostic-only.

\subsection*{Why the paper is theorem-led conditional rather than hybrid-led}

The proof-carrying core of the manuscript is already visible in
Sections~\ref{sec:primitive-roles}--\ref{sec:conditional-lift}. Those sections
define the theorem objects, state the internal theorem ladder, and culminate in
the conditional arithmetic lift. The vendor-backed material does not bear
theorem-level weight in that route. It supports primitive-role discipline, scope
discipline, and the refusal to collapse richer interaction structure into a
single proxy, but it does not discharge any theorem in
Sections~\ref{sec:fixed-package}--\ref{sec:conditional-lift}.

For that reason the paper is theorem-led conditional rather than hybrid-led. A
hybrid-led presentation would misplace the proof burden by making the empirical
material constitutive of the central theorem package. That is not the structure
of the manuscript. The theorem package stands on its internal proofs together
with the explicit external assumptions in the conditional arithmetic lift.

The paper is also conditional rather than unconditional for a precise reason:
the remaining arithmetic boundary has been isolated into an explicit external
contract. This makes the remaining literature dependence finite, named, and
visible at the theorem surface.

\subsection*{Out-of-scope claims}

Several stronger claims are intentionally excluded.

The paper does not prove an unconditional global canonicality theorem. It does
not show that every efficient operator in every possible regime collapses to one
universal representative family. It does not establish invariance under
arbitrary benchmark or cost perturbations. It does not use vendor-backed
interaction evidence as theorem support. And it does not let \(P3\) route
sensitivity substitute for genuine package change.

The paper also does not claim that the full Six-Birds or PICA interaction space
has been reduced to one theorem-bearing summary. The theorem ladder proven here
lives on a deliberately frozen slice. That restriction is part of the design of
the argument. The broader interaction space remains real, and the supporting
vendor-backed evidence is included precisely to keep that fact visible without
letting it carry theorem-level force.

Nor does the paper claim that the toy falsification phase supports the main
theorem. Its role was negative and disciplinary: it ruled out naive frontier
rhetoric in small finite spaces. The proof-bearing route remains the internal
Lean-backed theorem ladder plus the explicit external arithmetic contract.

\subsection*{What the theorem does license}

The main theorem of the paper is therefore the conditional theorem of Section~\ref{sec:conditional-lift}, not an unrestricted canonicality claim.

With the preceding exclusions in place, the positive scope of the paper is
clear. The manuscript proves that fixed packages saturate, that genuine
persistent growth requires package change, and that on a frozen evaluation
regime efficiency transfers across cone-equivalent representatives and yields
restricted alignment on the relevant support cone. It then lifts that result
conditionally into an arithmetic setting under an explicit external contract.
This scope is sufficient for the theorem-led foundations argument developed
here.

The paper also licenses a disciplined reading of the Six-Birds framework in
foundations. It shows that one can specialize the primitive vocabulary without
reducing the framework to a loose heuristic. In the present setting, the primitive roles
are asymmetric for principled reasons, and the theorem package depends on
respecting that asymmetry.

The manuscript should therefore be read as a theorem paper with an explicit
perimeter: a complete internal theorem ladder, a conditional arithmetic summit,
and a supporting empirical perimeter that constrains scope without carrying
proof.
